\section{Related Work}
\noindent\textbf{Vision-Language-Action Models.}
Vision-Language-Action (VLA) models build upon pretrained Vision-Language Models (VLMs) to enable language-conditioned robot control. Early work such as RT-2~\cite{rt22023arxiv} fine-tunes a VLM to output discretized action tokens. Subsequent efforts scale this paradigm through larger datasets and diverse embodiments~\cite{open_x_embodiment_rt_x_2023, octo_2023, khazatsky2024droid, kim2024openvla}, and more recent models have converged on a common architectural pattern: pairing a VLM backbone with a diffusion or flow-matching action head~\cite{li2024cogact, pi0, pi05, gr00tn1_2025, molmobot}. Despite rapid progress, recent stress tests reveal that these models remain brittle along two complementary axes: \emph{spatial generalization}, where success rates collapse under changes to object positions~\cite{zhou2025liberopro, fei25libero-plus}, and \emph{task generalization}, where models replay memorized trajectories rather than respond to changes in the instruction~\cite{zhou2025liberopro}. Closing these gaps has emerged as a central challenge in the development of robust VLAs.

\noindent\textbf{Spatial Grounding for Vision-Language-Action Models.}
A growing line of work addresses the mentioned generalization gaps by endowing the policy with explicit grounding signals indicating where the task-relevant object or region lies. One family relies on external grounding modules---typically VLMs or detectors that predict pixel coordinates~\cite{yuan2024robopoint}, 2D trajectories~\cite{li2025hamster}, or segmentation masks~\cite{roboground}---and convey their predictions to the policy either through language prompt~\cite{molmobot} or visual prompt~\cite{moka, wang2026vp, li2025hamster}. Another family instead trains the VLA itself to emit such spatial outputs as an intermediate chain-of-thought prediction before producing actions~\cite{coavla, molmoact2025, ecot, steerablepolicy}, with some further routing these predictions through dedicated injection modules to the action head~\cite{coavla}. Despite this diversity in source and routing, the spatial signal in prior work remains a 2D quantity tied to the image plane, and how this signal is represented and injected into the policy has received little systematic attention. We show that this is the true bottleneck, and that a simple change to both---lifting the signal into 3D and injecting it directly into the action head---unlocks substantial gains.

\noindent\textbf{3D-Enhanced Vision-Language-Action Models.}
The value of 3D information for manipulation has been demonstrated by a line of work that introduces explicit 3D representations into VLAs, ranging from point cloud alignment in 3D-based LLMs~\cite{zhen20243dvla} and ego-centric position encodings in the backbone~\cite{spatialvla} to dense spatial tokens injected into the action head~\cite{pointvla, falcon}. These methods consistently show that richer 3D information improves manipulation, but obtain it through dense scene-level representations that require a dedicated 3D encoder---and in some cases an additional foundation model---on top of the VLA. Our work pushes this observation to its minimal form: starting from a 2D grounding signal already available from off-the-shelf detectors or VLMs, a single task-relevant point lifted into 3D---encoding only the gripper-to-target displacement---already delivers substantial gains, with no additional 3D encoder and no changes to the backbone.